\documentclass[a4paper,12pt]{article}
\usepackage[utf8]{inputenc}
\usepackage[T1]{fontenc}
\usepackage[spanish]{babel}
\usepackage{geometry}
\usepackage{amsmath,amssymb}
\usepackage{multicol}
\usepackage{palatino}
\usepackage{mathpazo}
\usepackage{enumitem}
\geometry{margin=1.5cm}

\begin{document}
\thispagestyle{empty}
\large
%===================== EVALUACIÓN 1 =====================
\textit{Escuela 4‑117 “Ejército de los Andes| Evaluación: Principio de Pascal| 5º 3ª }

\noindent\textit{Apellido y Nombre:} \underline{\hspace{12cm}}
\begin{enumerate}[itemsep=0pt,parsep=0pt]
  \item Calcular la presión absoluta y la presión manométrica en el agua a $h=10{\,}\mathrm m$ (densidad $\rho=1000\:\mathrm{kg/m^3}$, $g=9{,}81\:\mathrm{m/s^2}$).
  \item ¿A qué profundidad en agua se alcanza $p_\text{abs}=1{,}5\:\mathrm{atm}$? (1 atm = $1{,}013\times10^5\:\mathrm{Pa}$).
  \item En una cañería se mide $h_\text{mmHg}=40\:\mathrm{mm}$ con manómetro de mercurio ($\rho_{\mathrm{Hg}}=13.6\times10^3\:\mathrm{kg/m^3}$). Calcular la presión absoluta en el fluido.
  \item Diferencia de presión en aceite ($\rho=850\:\mathrm{kg/m^3}$) entre dos puntos separados verticalmente $h_1=0{,}15\:\mathrm m$ y $h_2=0{,}25\:\mathrm m$.
  \item En una prensa hidráulica, el émbolo pequeño tiene radio $r_2=0{,}10\:\mathrm m$ y el grande $r_1=0{,}50\:\mathrm m$. ¿Qué fuerza $F_2$ debe aplicarse en el pequeño para elevar un automóvil de $m=1200\:\mathrm{kg}$?
  \item Si $A_1=20\:\mathrm{cm^2}$, $A_2=200\:\mathrm{cm^2}$ y $F_1=50\:\mathrm N$, hallar la fuerza de salida $F_2$.
\end{enumerate}

%===================== EVALUACIÓN 2 =====================
\textit{Escuela 4‑117 “Ejército de los Andes| Evaluación: Principio de Pascal| 5º 3ª}

\noindent\textit{Apellido y Nombre:} \underline{\hspace{12cm}}
\begin{enumerate}
  \item Calcular la presión absoluta en agua a la profundidad del Titanic: $h=3843\:\mathrm m$.
  \item ¿Qué profundidad en agua corresponde a $p_\text{man}=2{,}5\:\mathrm{kg/cm^2}$? (1 kgf/cm² = $9.81\times10^4\:\mathrm{Pa}$).
  \item Repetir el problema del manómetro de $h_\text{mmHg}=33\:\mathrm{mm}$ del práctico.
  \item En un circuito de aceite ($\rho=800\:\mathrm{kg/m^3}$) los desniveles son $h_1=0{,}10\:\mathrm m$, $h_2=0{,}20\:\mathrm m$, $h_3=0{,}32\:\mathrm m$. Hallar $\Delta p$ entre ambos extremos.
  \item Elevador hidráulico: radio del plato grande $r_1=0{,}50\:\mathrm m$, radio del chico $r_2=0{,}08\:\mathrm m$, $m=1000\:\mathrm{kg}$. Calcular $F_2$.
  \item ¿Qué fuerza $F$ requiere un pistón de área $A=10\:\mathrm{cm^2}$ para elevar agua hasta $h=12\:\mathrm m$?
\end{enumerate}

%===================== EVALUACIÓN 3 =====================
\textit{Escuela 4‑117 “Ejército de los Andes| Evaluación: Principio de Pascal| 5º 3ª}

\noindent\textit{Apellido y Nombre:} \underline{\hspace{12cm}}
\begin{enumerate}
  \item Presión manométrica en agua a $h=7{,}6\:\mathrm m$ y absoluta a $h=3843\:\mathrm m$.
  \item Profundidad en agua necesaria para $p_\text{abs}=3\times10^5\:\mathrm{Pa}$.
  \item En una prensa, discos de $r_1=2{,}00\:\mathrm m$ y $r_2=0{,}50\:\mathrm m$ elevan $m=1000\:\mathrm{kg}$. Determinar $F_2$.
  \item Área $A_1=0{,}001\:\mathrm{m^2}$, $A_2=0{,}10\:\mathrm{m^2}$, $F_1=100\:\mathrm N$. Hallar $F_2$.
  \item En una prensa de radios relacionados $r_1=3\,r_2$, se coloca $m_2=6\:\mathrm{kg}$ sobre el menor. ¿Qué masa $m_1$ puede levantarse en el mayor?
  \item Mismo caso, pero $m_2=16\:\mathrm{kg}$ y $r_1=2\,r_2$. Calcular $m_1$.
\end{enumerate}



\end{document}
